% Generated draft: requires mathematical and source review.
\documentclass[11pt]{article}
\usepackage[margin=25mm]{geometry}
\usepackage{fontspec}
\setmainfont{Latin Modern Roman}
\usepackage{amsmath,amssymb,mathtools,unicode-math}
\setmathfont{Latin Modern Math}
\usepackage{enumitem,needspace,xurl,hyperref}
\setlength{\parindent}{0pt}
\setlength{\parskip}{5pt}
\setlength{\emergencystretch}{3em}
\newcommand{\N}{\mathbb{N}}
\newcommand{\Z}{\mathbb{Z}}
\newcommand{\Q}{\mathbb{Q}}
\newcommand{\R}{\mathbb{R}}
\newcommand{\C}{\mathbb{C}}
\newcommand{\F}{\mathbb{F}}
\newcommand{\A}{\mathbb{A}}
\newcommand{\PP}{\mathbb{P}}
\newcommand{\E}{\mathbb{E}}
\newcommand{\eps}{\varepsilon}
\newcommand{\dd}{\,\mathrm{d}}
\begin{document}
\section*{UnsolvedMath ID: 20000216}
\textbf{Source problem code:} catalogue entry
\subsection{Definitions and mathematical statement}
The supplied verification field gives the evaluation map
\[
 \pi:\mathcal M_{g,n}(\mathbb P^r,\beta)\longrightarrow
 \operatorname{Conf}(\mathbb P^r,n),\qquad
 (f:C\to\mathbb P^r,p_1,\ldots,p_n)\longmapsto(f(p_1),\ldots,f(p_n)).
\]
The configuration target requires the open locus where marked images are distinct. The precise moduli component, map/embedding convention, and ground field are not reproduced. The verification field notes an initial source typo $\mathbb P^n$ where this displayed setup consistently uses $\mathbb P^r$.

\textbf{AIM degeneration problem 1.2.} Find understandable examples exhibiting behavior different from the expected one, especially unexpected tangent dimensions relevant to deciding dominance of $\pi$. The source expressly presents this as example hunting, not a request for a universal classification. It suggests using degeneration to propagate or analyze tangent-dimension defects; the exact hypotheses for any such deduction must be stated.

Original Remark 1.3 mentions attaching canonical curves and assumes nonnegative Brill--Noether number, conventionally $\rho(g,r,d)=g-(r+1)(g-d+r)$. It says dominance is known to fail in the indicated canonical examples and that a special canonical-type divisor gives a larger defect. Its divisor typography and final ``P3.12'' fragment are garbled; the report below interprets the relevant divisor as bicanonical, but the remark cannot be quoted as a fully recovered formula.

\textbf{Supplied research-aid context, not a verified solution.} Let $C\subset\mathbb P^3$ be a general smooth canonical genus-four $(2,3)$ complete intersection, and let $D$ be a general reduced divisor of degree $n\ge7$. For embedded deformations its normal bundle is $N_C=\mathcal O_C(2)\oplus\mathcal O_C(3)$, and the report claims the evaluation-kernel dimension
\[
 h^0(N_C(-D))=(9-n)_++(15-n)_+,
 \qquad(x)_+=\max(x,0).
\]
The Euler-expected value is $\max(0,24-2n)$. Thus it identifies degrees $10,11,12,13,14$ as failures of normal-bundle interpolation.

For $n=10,11,12$, it reports differential cokernel dimension $n-9$, equal to the codimension of the locus of configurations lying on a quadric, giving nondominance even though source/target dimension permits dominance. The quadric-incidence argument is important: in positive characteristic, differential nonsurjectivity alone need not exclude an inseparable dominant map. Field and genericity assumptions are not specified in the attempt.

At degree $12$, the report contrasts a general divisor with a reduced $D\in|2K_C|$. It claims
\[
 N_C(-D)\cong\mathcal O_C\oplus K_C,
 \qquad(h^0,h^1)=(5,5),
\]
whereas a general degree-twelve divisor gives $(3,3)$. This is a jump from the already defective generic value, not from the Euler-expected zero. These precise examples are unreviewed partial claims and do not by themselves extend a dominance theorem to every degeneration.
\subsection{Short English statement}
Find explicit tangent-dimension defects and use degeneration to study evaluation-map dominance. The attached genus-four example supplies a normal-bundle formula, quadric-incidence coranks, and a further bicanonical jump under missing field/genericity hypotheses.
\subsection{Sources}
\begin{itemize}
\item Local UnsolvedMath registry, permanent ID 20000216, source code catalogue entry; no proposer or proposal year supplied.
\item AIM workshop \emph{Degenerations in Algebraic Geometry} (2016), problem-session document, Problem 1.2 and original Remark 1.3; supplied locator \url{https://aimath.org/pastworkshops/degenalggeomproblems.pdf}. Canonical locator: aim-algebraic-geometry-notes.json, notes[215], tag ``problem.'' The original literature field is empty. No PDF was fetched or generated here.
\item Research attempt 1: partial\_result, status at run ``partially\_solved,'' machine-assigned rigor ``proved,'' novelty confidence ``low,'' with an explicit machine-generated/not-peer-reviewed/verify-before-citing caveat. No separate primary reference for the exact tangent calculation is supplied.
\end{itemize}
\end{document}
